\documentclass[10pt,a4paper]{amsart}
\usepackage[margin=19mm]{geometry}
\usepackage{amsmath,amssymb,mathtools}
\usepackage[scr=rsfs,cal=euler]{mathalfa}
\usepackage{xfrac}
\usepackage[unicode,hidelinks,bookmarks=false]{hyperref}
\newcommand{\ie}{i.e.\ }
\newcommand{\cf}{cf.\ }
\newcommand{\resp}{respectively\ }
\newcommand{\Subsection}{\subsection}
\providecommand{\nobreakdash}{\nobreak\hbox{-}}
\setlength{\emergencystretch}{2em}
\multlinegap0pt
\usepackage{bm}
\usepackage{bbm}
\usepackage{dsfont}
\usepackage{xspace}
\usepackage{cancel}
\usepackage{mathtools}
\usepackage[shortlabels]{enumitem}
\usepackage{amsthm}
\makeatletter
\@ifundefined{theorem}{\newtheorem{theorem}{Theorem}}{}
\@ifundefined{lemma}{\newtheorem{lemma}[theorem]{Lemma}}{}
\@ifundefined{proposition}{\newtheorem{proposition}[theorem]{Proposition}}{}
\makeatother
\newcommand{\D}{\mathfrak{D}}
\newcommand{\Mod}[1]{\ (\mathrm{mod}\ #1)} %For mod in paranthesis
\newcommand{\Q}{\mathbb{Q}}
\newcommand{\Z}{\mathbb{Z}}
\newcommand{\calC}{\mathcal{C}}
\newcommand{\calE}{\mathcal{E}}
\renewcommand{\a}{\alpha}
\title[Arithmetic Aspects of Number Fields Generated by Polynomial ]{Arithmetic Aspects of Number Fields Generated by Polynomial Families}
\author{Rupam Barman, Anuj Jakhar, Ravi Kalwaniya,, Prabhakar Yadav}
\date{Source: arXiv:2602.19726v1 (CC BY 4.0)}
\begin{document}
\maketitle

\noindent\emph{Source and licence.} Arithmetic Aspects of Number Fields Generated by Polynomial Families. Version arXiv:2602.19726v1 (CC BY 4.0), distributed under the Creative Commons Attribution 4.0 International licence, as recorded in the arXiv licence metadata for this exact version and shown on the abstract page https://arxiv.org/abs/2602.19726v1. Full rights notice: licenses/arxiv-2602.19726-CC-BY-4.0.txt.\par\smallskip

\noindent\textit{Editorial scope.} Counted unit: the Lemma whose source label is \texttt{dedekind} in the article (source0.tex lines 774--781, source label \texttt{dedekind}), delivered together with the article's own complete original proof of it (source0.tex lines 964--1172, one \texttt{proof} environment of 209 lines). No mathematics is written, repaired or re-derived: statement and proof are the article's own text line for line. The proof is detached from the statement in the source: the article states the result at lines 774--781 and prints its proof at lines 964--1172, and the proof environment's own header names the statement, so the association is the article's own. Required context retained uncounted: disc of f(x) (lines 533--542); monogeneity of f(x) (lines 549--567); binomial expansion (lines 768--771); lemma:3.6 (lines 784--786); 2.06 (lines 822--829); 0.009 (lines 891--894); 4.4:prop (lines 919--925). Every definition, display and label of those lines is retained unchanged. Omitted from this package and not redistributed: 1--532, 543--548, 568--767, 772--773, 782--783, 787--821, 830--890, 895--918, 926--963, 1173--1803. Nothing inside a retained excerpt is omitted or altered: every retained body is delivered exactly as the article's lines are written, so no omission receipt of any kind accompanies this package. Citations: the retained excerpts carry no citation command at all, so no bibliography entry of the article is transcribed and no reference is redistributed. Reference adaptation: none was needed and none was made; every reference inside the delivered unit resolves inside it, so no unresolved reference or citation is produced. Printed slips kept literal: no printed word, symbol, hypothesis or number of the counted statement or of its proof was altered. Layout and class: the article's own class is \texttt{amsart} with packages including float, geometry, amscd, amsmath, amsxtra, amsthm, amssymb, stmaryrd, xr, mathrsfs, mathtools, enumerate, commath, comment; the wrapper is the lane's pinned \texttt{amsart} editorial wrapper at 10pt on A4, which restates the theorem-like environments the retained text needs and adds the article's own macro definitions that the retained text needs (\texttt{\textbackslash D}, \texttt{\textbackslash Mod}, \texttt{\textbackslash Q}, \texttt{\textbackslash Z}, \texttt{\textbackslash a}, \texttt{\textbackslash calC}, \texttt{\textbackslash calE}), with extra wrapper packages (bm, bbm, dsfont, xspace, cancel, mathtools, enumitem). The delivered numbering is the unit's own. Assets: no retained span contains a figure environment, a graphics-inclusion command, a PDF-page inclusion, a table or a TikZ picture; no image file is copied into this package, the package declares no asset dependency and no image file is redistributed with it. Licence: version arXiv:2602.19726v1 of this article is distributed under the Creative Commons Attribution 4.0 International licence, as recorded in the arXiv licence metadata for this exact version and shown on the abstract page https://arxiv.org/abs/2602.19726v1. A full rights notice is delivered at licenses/arxiv-2602.19726-CC-BY-4.0.txt. Characters: every retained excerpt is pure ASCII, so the delivered engine, XeLaTeX with its default Latin Modern text font, reports no missing character.
\begin{theorem} \label{disc of f(x)}
Let $a,c \in \mathbb{Z}$ with $c \neq 0$, and let $m,n,k \in \mathbb{N}$ with $km>n$. 
Suppose that the polynomial $f(x) = (x^k + c)^m - a x^n \in \mathbb{Z}[x]$ is irreducible over $\mathbb{Z}$ and let $\theta$ be a root of $f(x)$. Then %This form is accurate other form of writing will be more messy.
\begin{equation}
    \D_f
=(-1)^{\binom{km}{2}+(km+n+t)(m-1)}\,a^{\,k(m-1)}c^{m(n-1)}\,
\left[\,(km)^{\,k_1m}c^{k_1m-n_1} - a^{k_1}n^{\,n_1}\,(km-n)^{\,k_1m-n_1}\right]^{t}.
\end{equation}
where $n = n_1 t$, $k = k_1 t$ and $t=\gcd(n,k)$.
\end{theorem}

\begin{theorem} \label{monogeneity of f(x)}
	Let $K=\mathbb Q(\theta)$ be an algebraic number field with $\theta$ in the ring $\Z_K$ of algebraic integers of $K$ having  minimal polynomial   $f(x) = (x^k+c)^m-ax^n$ of degree $km$ over $\mathbb Q$ with   $km>n\geq{1}$. Let $\gcd(n,k) = t$ and $n = n_1t,\, k = k_1t.$ A prime divisor $p$ of the discriminant  $\D_f$ of $f(x)$ does not divide $[\Z_{K} : \Z[\theta]]$ if and only if $p$  satisfies one of the following conditions: 
	 \begin{enumerate}[label=\textup{(\roman{enumi})}]
    \item  If $p\mid c$  and $p\mid a$, then $m=1$ and $p^2\nmid c$;
    \item  If $p\mid c,~~p\nmid a$ and $n=1$  with $j\geq 1$ as the highest power of $p$ dividing $km-1$, then either $p\mid c_2$ and $p\nmid a_1$ or $p\nmid c_2[a^ka_1^{\,km-1}-(-mac_2)^{\,km-1}]$, where $c_2=\frac{c}{p}$ and $a_1=\frac{a^{p^j}-a}{p}$;
    \item  If $p\mid c,~~p\nmid a$ and $n\geq 2$ with $\ell\geq 0$ as the highest power of $p$ dividing $k-n$, then $m=1$ and  either
$p\mid a_{1}$ and $p\nmid c_{2}$, or $p\nmid a_{1}c_{2}^{\,n-1}
\bigl[a^{n_1}a_{1}^{\,k_1-n_1}-(-c_{2})^{\,k_1-n_1}\bigr]$, where $c_2=\frac{c}{p}$ and $a_1=\frac{a^{\,p^{\ell}}-a}{p}$;
    \item  If $p\nmid c$ and $p\mid a$ and further;
    \begin{itemize}
        \item[\textup{(a)}] If $m\geq 2$, then $p^2\nmid a$;
        \item [\textup{(b)}] If $m=1$ with $\ell\geq 1$ as the highest power of $p$ dividing $k$, then either $p\mid a_2$ and $p\nmid c_1$ or $p\nmid a_2\bigl[a_{2}^{k_{1}}(-c)^{n_{1}}-(-c_{1})^{k_{1}}\bigr]$, where $a_2=-\frac{a}{p}$ and $c_1=\frac{(-c)^{\,p^{\ell}}+c}{p}$;
    \end{itemize}
    \item If $p \nmid ca$ and $p \mid k$, where $k = p^{j}s'$ and $n = p^{j}s$ with 
$p \nmid \gcd(s,s')$, then the polynomials $\frac{1}{p}[a_1\,x^{n} - p m\, t(x)(x^{k} + c)^{\,m-1}]$ and $(x^{s'} + c)^m - a x^s$ are coprime, where $a_1=\frac{a^{\,p^{j}}-a}{p}$ and\\ $t(x) = \frac{1}{p} \left( \sum_{i=0}^{p^{j}-1} \binom{p^{j}}{i} x^{\,s' \,i} \,c^{\,p^{j}-i} - c \right)$;
     \item If $p\nmid cak$ and $p\mid m$, then $p \mid n$ and $p^2\nmid (a^{p}-a)$;
    \item  If $p\nmid cakm$, then $p^2\nmid (\calC-\calE).$
	\end{enumerate}
\end{theorem}    

\begin{lemma}\label{binomial expansion} { Let $k\geq 1$ be the highest power of a prime $p$ dividing a number $n=p^{k}s'$ and $c$ be an integer not divisible by $p$. If $\bar{g}_1(x)\cdots \bar{g}_r(x)$  is the factorization of $x^{s'}-\bar{c}$ into  a product of distinct irreducible polynomials over $\Z/p\Z$ with $g_i(x)\in\mathbb Z[x] $ monic, then 
		$$x^n-c=(g_1(x)\cdots g_r(x)+pH_1(x))^{p^k}+pg_1(x)\cdots g_r(x)H_2(x)+p^2H_3(x)+c^{p^k}-c$$ 
		\noindent for some polynomials  $H_1(x), H_2(x), H_3(x)\in \mathbb Z[x]$.}
\end{lemma}

\begin{lemma}\label{lemma:3.6}
	Let $f(x)=(x^k+c)^m-ax^n\in\Z[x]$ be a polynomial of degree $km$. Let $p$ be a prime number such that $n=p^js,~k=p^js',$ where $p\nmid\gcd(s,s')$. Define $h(x)=(x^{\,s'}+c)^m-ax^s\in\Z[x]$. Then $f(x)$ can be expressed as \[f(x) =  h(x)^{\,p^{j}}+(a^{\,p^{j}} + a)\,x^{n} - p m\, t(x)(x^{k} + c)^{\,m-1}+ p\,h(x)\,h_{1}(x) + p^{2}\,h_{4}(x),\] for some $h_1(x),\, h_4(x)\in\Z[x]$ and $t(x) = \frac{1}{p} \left( \sum_{i=1}^{p^{j}} \binom{p^{j}}{i} x^{\,s' p^{j}} c^{\,p^{j}-i} - c \right) \in \mathbb{Z}[x].$
\end{lemma}

\begin{proposition}\label{2.06}
	Let $f(x), ~n_1, ~k_1, ~\calC$ and $\calE$ be as defined in Theorem \ref{monogeneity of f(x)}. Let $u_1, u_2\in \Z$ satisfy $ ku_1+nu_2=t$ and $p$ be a prime such that $p\nmid ackm(km-n)$. Furthermore, suppose that there exist integers $ \alpha_1$ and $\alpha_2$ such that $\alpha_1 \equiv \frac{nc}{km-n} \pmod{p^2}$ and $\alpha_2\equiv a^{-1} \left(\frac{kmc}{km-n}\right)^m \pmod{p^2}$. Then the following are equivalent for $i=1,2$.
	\begin{itemize}
		\item[\textup{(i)}] $p^i\mid (\calC-\calE)$;
		\item [\textup{(ii)}] $\alpha_2^{k_1}\equiv\alpha_1^{n_1}\pmod{p^i}$;
		\item [\textup{(iii)}] If $\beta$ satisfies $\beta^t \equiv \alpha_1^{u_1} \alpha_2^{u_2} \pmod{p^i}$, then $\beta$ is a repeated root of $f(x)$ modulo $p^i$.
	\end{itemize}
\end{proposition}

		\begin{align}\label{0.009}
			\beta^k &\equiv \frac{nc}{km-n}\pmod{p},\nonumber\\
			&\equiv \alpha_1 \pmod{p}.
		\end{align}

\begin{proposition} \label{4.4:prop}
    Assume the notation introduced in Proposition \ref{2.06}. Let $\a$ be an integer such that $\a \equiv \a_1^{\,u_1}\a_2^{\,u_2} \Mod{p^2}$. Let $p$ be the prime such that $p \mid (\calC - \calE)$ and $p \nmid ackm(km-n)$. Then the following are equivalent:
    \begin{itemize}
        \item[\textup{(i)}] $(\alpha^{\,k_1} + c)^m - a \alpha^{\,n_1} \equiv 0 \Mod{p^2}$;
        \item[\textup{(ii)}] $\alpha_1^{\,n_1} \equiv \alpha_2^{\,k_1} \Mod{p^2}$.
    \end{itemize}
\end{proposition}

\begin{lemma}\label{dedekind}
	Let  $f(x) \in \Z[x]$ be a monic irreducible polynomial  having factorization $\bar{g}_1(x)^{e_1} \cdots \bar{g}_{t}(x)^{e_{t}}$ modulo a prime $p$ as a product of powers of distinct irreducible polynomials over $\Z/p\Z$ with $g_i(x) \in \Z[x]$ monic. Let $K=\Q(\theta)$ with $\theta$ a root of $f(x)$.   Then the following statements are equivalent:
	\begin{itemize}
		\item[\textup{(i)}] $p$ does not divide $[\Z_K:\Z[\theta]]$.
		\item[\textup{(ii)}] For each $i$, we have $e_i =1 $ or $\overline g_i (x)$ does not divide $\overline M(x)$, where $M(x) = \frac{1}{p}(f(x) -  g_1 (x)^{e_1} \cdots  g_{t} (x)^{e_{t}} )$.
		\item[\textup{(iii)}] $f(x)$ does not belong to the ideal $\langle p, g_i(x)\rangle^2$ in $\Z[x]$ for any $i$, $1\leq i\leq t$.
	\end{itemize}
\end{lemma} 

\begin{proof}[Proof of Theorem~\ref{monogeneity of f(x)}]
Let $p$ be a prime that divides $\D_f$. In view of Lemma~\ref{dedekind}, we have $p \nmid [\mathbb{Z}_K : \mathbb{Z}[\theta]]$
if and only if $f(x) \not\in \langle p,\, g(x) \rangle^2$
for every monic polynomial $g(x) \in \mathbb{Z}[x]$ that is irreducible modulo $p$. Moreover, note that
$f(x) \not\in \langle p, g(x) \rangle^2$
whenever $\overline{g}(x)$ is not a repeated factor of $\bar{f}(x)$.\\[2mm]
\textbf{Case (i).} Suppose $p \mid c$ and $p \mid a$. In this case, $f(x) = (x^k + c)^m - a x^n \equiv x^{km} \pmod{p}.$ It is easily checked that $f(x) \in \langle p,\, x \rangle^2 $ if and only if $p^2 \mid c^m.$
Therefore, by Lemma~\ref{dedekind}, $p \nmid [\mathbb{Z}_K : \mathbb{Z}[\theta]]
\textit{ if and only if }
f(x) \notin \langle p,\, x \rangle^2 $ if and only if $p^2 \nmid c^m.$ Since $p \mid c$, condition $p^2 \nmid c^m$ can hold only when $m=1$ and $p^2 \nmid c$. This completes the proof in this case.\\[2mm]
\textbf{Case (ii).} Suppose that $p\mid c$, $p\nmid a$ and $n=1$.
Since $p\mid \D_f$, $p\mid c$ and $p\nmid a$, it follows from
\eqref{disc of f(x)} that $p\mid (km-1).$ Write $km-1=p^{j}s,
\, p\nmid s.$ In this situation, $f(x)\equiv x\bigl(x^{km-1}-a\bigr)
      \equiv x\bigl(x^{s}-a\bigr)^{p^{j}}
      \pmod{p}.$
Set $h(x)=x^{s}-a.$ Raising both sides to the power $p^{j}$ and using the Binomial theorem, we obtain the following.
\begin{equation}\label{eq:4.9}
x^{km-1}
= h(x)^{p^{j}} + a^{p^{j}} + p\,h(x)h_{1}(x),
\end{equation}
for some polynomial $h_{1}(x)\in\mathbb{Z}[x]$. Since $p\mid c$, we may expand $f(x)$ as
\begin{align}\label{eq:4.10}
f(x)
&= x^{km}-ax+\sum_{i=1}^{m-1}\binom{m}{i}x^{ki}c^{m-i} \nonumber \\
&= x(x^{km-1}-a)+m x^{k(m-1)}c + p^{2}h_{2}(x),
\end{align}
for some $h_{2}(x)\in\mathbb{Z}[x]$. Write $h(x)=g_{1}(x)\cdots g_{t}(x)+pH(x),$ where $g_{1}(x),\ldots,g_{t}(x)$ are distinct monic polynomials that are irreducible modulo $p$ and $H(x)\in\mathbb{Z}[x]$.
Substituting this expression into \eqref{eq:4.10} and using
\eqref{eq:4.9}, we obtain
\begin{align*}
f(x)
&= x\,h(x)^{p^{j}} + m c x^{k(m-1)} + (a^{p^{j}}-a)x
   + p\,h(x)h_{3}(x) + p^{2}h_{4}(x) \\
&= x\left(\prod_{i=1}^{t} g_i(x)+pH(x)\right)^{p^{j}}
   + x\bigl(m c x^{k(m-1)-1} + (a^{p^{j}}-a)\bigr)
   + p\,h(x)h_{3}(x) + p^{2}h_{4}(x),
\end{align*}
for some $h_{3}(x),\,h_{4}(x)\in\mathbb{Z}[x]$. Observe that $x,\,\bar g_{1}(x),\ldots,\bar g_{t}(x)$ are pairwise distinct irreducible factors of $\overline{f}(x)$.
Since $j\ge1$, it follows that $f(x)\in\langle p,\, g_i(x)\rangle^{2}$
for some $i$ if and only if the term $m c x^{k(m-1)-1} + (a^{p^{j}}-a)$ lies in $\langle p,\, g_i(x)\rangle^{2}$. Write $a^{p^{j}}-a=pa_{1},\,c=pc_{2}.$ Since $p\mid (km-1)$ implies that $p\nmid m$. Then $f(x)\in\langle p,\, g_i(x)\rangle^{2}$ for some $i$ if and only if either $p\mid c_{2}$ and $p\mid a_{1}$, or $p\nmid c_{2}$ and the polynomials $\bar m\,\bar c_{2}\,x^{k(m-1)-1}+\bar a_{1}
\quad\text{and}\quad
x^{s}-\bar a$ have a common root modulo $p$. Let $\eta$ be such a common root. Then $m c_{2}\eta^{k(m-1)-1}=-a_{1},
\,
\eta^{s}=a.$ Eliminating $\eta$, we obtain the common-root condition $a^ka_2^{km-1}=(mc)^{km-1}.$ Therefore, we conclude that $p$ does not divide $[\Z_k :\Z[\theta]]$ if and only if either $p\mid c_2$ and $p\nmid a_1$ or $p$ does not divide $c_2[a^ka_1^{\,km-1}-(-mac_2)^{\,km-1}]$. This completes the proof in the present case. \\[2mm]
\textbf{Case (iii).} Suppose that $p\mid c$ and $p\nmid a$. In this case,
$f(x)\equiv x^{n}\bigl(x^{km-n}-a\bigr)\pmod{p}.$ Since $n\ge2$, it follows that $x$ is a repeated root of $f(x)$ modulo $p$. Arguing as in Case~(i), we obtain $f(x)\notin\langle p,x\rangle^2
\quad \Longleftrightarrow \quad
p^{2}\nmid c^{m}.$ As $p\mid c$, this is possible only when $m=1$ and $p^{2}\nmid c$. Hence, we may assume throughout this case that $m=1$. Then $f(x)=x^{k}-ax^{n}+c$ and modulo $p$ we have $f(x)\equiv x^{n}\bigl(x^{k-n}-a\bigr).$ Let $\ell\ge0$ be the largest integer such that $p^{\ell}\mid(k-n)$.
\medskip
\emph{First, assume that $\ell=0$, i.e., \ $p\nmid(k-n)$.}
Since $p\mid c$ and $p\nmid a$, the reduction $\overline{f}(x)=x^{n}\bigl(x^{k-n}-\bar a\bigr)$ has $x$ as its only possible repeated irreducible factor modulo $p$. In this situation, it is easy to check that $f(x)\in\langle p,x\rangle^{2}
\textit{ if and only if }
p^{2}\mid c.$ Thus, when $\ell=0$, the desired conclusion follows immediately.\\
Now assume that $\ell\ge1$ and write $k-n=p^{\ell}s,
\quad p\nmid s.$ Set $h(x)=x^{s}-a.$ We may write $h(x)=g_{1}(x)\cdots g_{t}(x)+pH(x),$ where $g_{1}(x),\dots,g_{t}(x)$ are distinct monic polynomials that are irreducible modulo $p$ and $H(x)\in\mathbb{Z}[x]$.
Applying Lemma~\ref{binomial expansion} to $h(x)$, we obtain
\[
f(x)=x^{n}\Biggl[
\left(\prod_{i=1}^{t} g_i(x)+pH(x)\right)^{p^{\ell}}
+pT(x)\prod_{i=1}^{t} g_i(x)
+p^{2}U(x)
-a+a^{p^{\ell}}
\Biggr]+c,
\]
for some polynomials $T(x),U(x)\in\mathbb{Z}[x]$. Observe that $x,\bar g_{1}(x),\ldots,\bar g_{t}(x)$ are pairwise distinct
irreducible factors of $\overline{f}(x)$.
Since $\ell\ge1$, the first three terms inside the brackets are in
$\langle p,g_i(x)\rangle^{2}$ for each $1\le i\le t$.
Consequently, $f(x)\in\langle p,g_i(x)\rangle^{2}$ for some $i$ if and only if the remaining term $\bigl(a^{p^{\ell}}-a\bigr)x^{n}+c$
lies in $\langle p,\, g_i(x)\rangle^{2}$. Write $a^{p^{\ell}}-a=pa_{1},
\, c=pc_{2}.$
Then $f(x)\in\langle p,\,g_i(x)\rangle^{2}$ for some $i$ if and only if
either $p\mid a_{1}$ and $p\mid c_{2}$, or $p\nmid a_{1}$ and the
polynomials $\bar a_{1}x^{n}+\bar c_{2}
\quad\text{and}\quad
x^{k-n}-\bar a$ have a common root modulo $p$. Let $\gamma$ be such a common root. Then
\[
a_{1}\gamma^{n}+c_{2}=0,
\qquad
\gamma^{k}=a\gamma^{n}.
\]
Simplifying these relations and using $k=k_{1}t$ and $n=n_{1}t$, we obtain
\[
(-c_{2})^{\,k_{1}-n_{1}}=a^{\,n_{1}}a_{1}^{\,k_{1}-n_{1}}.
\]
Hence, the above polynomials have a common root modulo $p$ if and only if
\[
(-c_{2})^{\,k-n}=a^{\,n}a_{1}^{\,k-n}.
\]
By Lemma~\ref{dedekind}, we conclude that
$p\nmid[\mathbb{Z}_{K}:\mathbb{Z}[\theta]]$ if and only if either
$p\mid a_{1}$ and $p\nmid c_{2}$, or
$
p\nmid a_{1}c_2
\bigl[a^{n_{1}}a_{1}^{\,k_{1}-n_{1}}-(-c_{2})^{\,k_{1}-n_{1}}\bigr].$
This completes the proof in the present case.\\[2mm]
\textbf{Case (iv).} Suppose that $p\nmid c$ and $p\mid a$. Then $f(x) \equiv (x^{k}+c)^{m} \pmod{p}.$ There are two possibilities to consider, namely $m\ge 2$ and $m=1$.\\
\emph{First, assume that $m\ge 2$.}
Let $g(x)$ be an irreducible factor of $x^{k}+c$ over $\mathbb{Z}/p\mathbb{Z}$.
Then $f(x)\in\langle p,g(x)\rangle^{2}$ if and only if $p^{2}\mid a$.
Hence, by Lemma~\ref{dedekind}, we obtain $p\nmid [\mathbb{Z}_{K}:\mathbb{Z}[\theta]]
\quad \Longleftrightarrow \quad
p^{2}\nmid a.$\\
\emph{Now assume that $m=1$.}
In this case $f(x)=x^{k}-ax^{n}+c,$ and since $p\mid a$, we have $f(x)\equiv x^{k}+c \pmod{p}.$ Moreover, as $p\mid \D_f$, it follows from \eqref{disc of f(x)} that $p\mid k$. Write $k=p^{\ell}s \, \text{ with }\, p\nmid s.$ By the Binomial theorem, we then obtain $f(x)\equiv (x^{s}+c)^{\,p^{\ell}} \pmod{p}.$ Let $\bar g_{1}(x),\ldots,\bar g_{r}(x)$ be the factorization of
$x^{s}+\bar c$ over $\mathbb{Z}/p\mathbb{Z}$, where
$g_i(x)\in\mathbb{Z}[x]$ are monic, pairwise distinct and irreducible
modulo $p$. We may therefore write $x^{s}+c = g_{1}(x)\cdots g_{r}(x) + pH(x)$ for some polynomial $H(x)\in\mathbb{Z}[x]$. Set $h(x)=x^{s}+c$.
Then $f(x)=h(x^{p^{\ell}})-ax^{n}.$ Applying Lemma~\ref{binomial expansion} to $h(x)$, we obtain
\begin{equation}\label{eq:case-iii-expansion}
f(x)= \left(\prod_{i=1}^{r} g_i(x) + pH(x)\right)^{p^{\ell}}
   + pT(x)\prod_{i=1}^{r} g_i(x)
   + p^{2}U(x) + c + (-c)^{p^{\ell}} - ax^{n},
\end{equation} 
for some polynomials $T(x),U(x)\in\mathbb{Z}[x]$.
Since $\ell\ge 1$, the first three terms on the right-hand side of
\eqref{eq:case-iii-expansion} lie in the ideal
$\langle p, g_i(x)\rangle^{2}$ for each $1\le i\le r$.
Therefore,
$f(x)\in\langle p, g_i(x)\rangle^{2}$ for some $i$
if and only if $-ax^{n}+c+(-c)^{p^{\ell}} = p(a_{2}x^{n}+c_{1})$
belongs to $\langle p, g_i(x)\rangle^{2}$. It follows that
$p(a_{2}x^{n}+c_{1})\in\langle p, g_i(x)\rangle^{2}$
for some $i$ if and only if either $p\mid a_{2}$ and $p\mid c_{1}$,
or $p\nmid a_{2}$ and the polynomials
$\overline{a}_{2}x^{n}+\overline{c}_{1}$ and $x^{s}+\overline{c}$
have a common root modulo $p$. Let $\gamma$ be such a common root.
Then $a_{2}\gamma^{n}+c_{1}=0
\quad \text{and} \quad
\gamma^{k}+c=0.$ Writing $k=k_{1}t$ and $n=n_{1}t$, we obtain
\[
a_{2}(\gamma^{t})^{n_{1}}=-c_{1}
\quad \text{and} \quad
(\gamma^{t})^{k_{1}}=-c.
\]
The rise of the first equation to the power $k_{1}$ and the second to the power $n_{1}$ yields the following equation.
\[
a_{2}^{k_{1}}(\gamma^{t})^{n_{1}k_{1}}
=(-c_{1})^{k_{1}}
\quad \text{and} \quad
(\gamma^{t})^{k_{1}n_{1}}
=(-c)^{n_{1}}.
\]
Comparing these expressions, we conclude that $a_{2}^{k_{1}}(-c)^{n_{1}} = (-c_{1})^{k_{1}}.$ Therefore, the polynomials
$\overline{a}_{2}x^{n}+\overline{c}_{1}$ and $x^{s}+\overline{c}$
are coprime modulo $p$ if and only if $p\nmid \bigl(a_{2}^{k_{1}}(-c)^{n_{1}}-(-c_{1})^{k_{1}}\bigr).$ Therefore, by Lemma~\ref{dedekind}, we obtain $p\nmid [\mathbb{Z}_{K}:\mathbb{Z}[\theta]]$ if and only either $p\mid a_2$ and $p\nmid c_1$ or $p\nmid a_2\bigl[a_{2}^{k_{1}}(-c)^{n_{1}}-(-c_{1})^{k_{1}}\bigr]$.\\[2mm]
\textbf{Case (v).} Suppose $p \nmid ca$ and $p \mid k$. By Theorem~\ref{disc of f(x)}, we have 
$p \mid (\mathcal{C} - \mathcal{E})$, and hence $p \mid n(km - n).$
Since $p \nmid a$ and $p \nmid c$, we obtain $p \mid n$. Write $k = p^{j}s' \quad \text{and} \quad n = p^{j}s,$ with $j \ge 1$ and $p \nmid \gcd(s,s')$. Then
\[
f(x) \equiv \Bigl( (x^{s'} + c)^m - a x^s \Bigr)^{p^{j}} \pmod{p}.
\]
 Denote  $ (x^{s'} + c)^m - a x^s$ by $h(x)$ so that $f(x)=h(x^{p^j})$. Then one can easily check that $f(x)\equiv h(x)^{p^j} \pmod{p}$. Write $h(x) = g^{e_1}_1(x)\cdots g^{e_d}_d(x) + ph_1(x)$, $e_i>0$, where $g_1(x), \ldots,$ $g_d(x)$ are monic polynomials which are distinct as well as irreducible modulo $p$ and $h_1(x) \in \Z[x]$. Using Lemma \ref{lemma:3.6}, we have
\begin{align*}
    f(x) &=  h(x)^{\,p^{j}}+(a^{\,p^{j}} - a)\,x^{n} - p m\, t(x)(x^{k} + c)^{\,m-1}+ p\,h(x)\,h_{1}(x) + p^{2}\,h_{4}(x)\\
    &=\left(\prod\limits_{i=1}^{d}g^{e_i}_i(x)\right)^{\,p^j}+(a^{\,p^{j}} - a)\,x^{n} - p m\, t(x)(x^{k} + c)^{\,m-1}+ p\,\left(\prod\limits_{i=1}^{d}g^{e_i}_i(x)\right)\,h_{1}(x) + p^{2}\,h_{4}(x)
\end{align*}
for some polynomials $h_1(x),\,h_4(x)\in\Z[x]$. Write $f(x)=\left(g_1(x)^{e_1} \cdots g_d(x)^{e_d}\right)^{p^j}+p M(x)$ for some $M(x)\in\Z[x].$ Since  $j>0$, by Lemma \ref{dedekind}, we see that $p$ does not divide $[\Z_K :\Z[\theta]]$ if and only if $\overline{M}(x)$ is coprime to $\overline{h}(x)$, which holds if and only if the polynomial $\frac{1}{p}[(a^{\,p^{j}} - a)\,x^{n} - p m\, t(x)(x^{k} + c)^{\,m-1}]$ is coprime to $h(x)$ modulo $p$. This proves the theorem in this case.\\[2mm]
\textbf{Case (vi).} Suppose $p \nmid cak$ and $p \mid m$. By Theorem~\ref{disc of f(x)}, we have $p \mid n$. Write $n = p^{\ell}s \quad \text{and} \quad m = p^{\ell}s',$ where $\ell \ge 1$ and $p \nmid \gcd(s,s')$. Define $h(x) = (x^{k}+c)^{s'} - a x^{s}.$ Then it is easy to check that $f(x) \equiv h(x)^{\,p^{\ell}} \pmod{p}.$
We first expand $(h(x)+a x^{s})^{p^{\ell}}$ using the binomial theorem:
\[
(h(x)+a x^{s})^{p^{\ell}} = h(x)^{p^{\ell}} + pH(x) + a^{p^{\ell}}x^{n},
\]
for some polynomial $H(x)\in\mathbb{Z}[x]$. Since $h(x)=(x^{k}+c)^{s'}-a x^{s}$, it follows that
\[
(x^{k}+c)^{m}
= h(x)^{p^{\ell}} + p\,h(x)H(x) + (a^{p^{\ell}}-a+a)x^{n}.
\]
Consequently, we may write the following.
\begin{equation}\label{eq:5.1}
f(x) = h(x)^{p^{\ell}} + p\,h(x)h_{1}(x) + (a^{p^{\ell}}-a)x^{n},
\end{equation}
for some polynomial $h_{1}(x)\in\mathbb{Z}[x]$. Next, write $h(x)=g_{1}(x)^{e_{1}}\cdots g_{d}(x)^{e_{d}} + p\,h_{2}(x),
\quad e_i>0,$ where $g_{1}(x),\ldots,g_{d}(x)$ are distinct monic polynomials that are irreducible modulo $p$ and $h_{2}(x)\in\mathbb{Z}[x]$. Substituting this expression into \eqref{eq:5.1}, we obtain
\begin{align*}
f(x)
&= \left(\prod_{i=1}^{d} g_{i}(x)^{e_{i}} + p\,h_{2}(x)\right)^{p^{\ell}}
   + p\,h(x)h_{1}(x) + (a^{p^{\ell}}-a)x^{n} \\
&= \left(\prod_{i=1}^{d} g_{i}(x)^{e_{i}}\right)^{p^{\ell}}
   + p\,h(x)h_{1}(x) + p^{2}h_{3}(x) + (a^{p^{\ell}}-a)x^{n},
\end{align*}
for some polynomial $h_{3}(x)\in\mathbb{Z}[x]$. Thus, we may write 
$
f(x)=\left(g_{1}(x)^{e_{1}}\cdots g_{d}(x)^{e_{d}}\right)^{p^{\ell}} + pM(x),
$ for some $M(x)\in\mathbb{Z}[x]$. Since $\ell>0$, Lemma~\ref{dedekind} implies that $p \nmid [\mathbb{Z}_{K}:\mathbb{Z}[\theta]]$ if and only if $\overline{M}(x)$ is co-prime to $\overline{h}(x)$ modulo $p$.
This holds if and only if the polynomial $a_1 x^{n}$ is co-prime to $h(x)$ modulo $p$,
which is equivalent to condition $p\nmid a_1$, where $a_1=\frac{a^{p^{\ell}}-a}{p}.$ This condition also simplifies to $p^{2}\nmid (a^{p}-a)$. This completes the proof in this case.\\[2mm] 
\textbf{Case (vii).} Now consider the last  case where $p \nmid ackm $. Furthermore, as $p\mid \D_f$ and $p\nmid km$ and $k_1m-n_1 \neq 0$, it follows that $p \nmid n(km-n)$. We claim that there exists an integer $\alpha$ such that the polynomial $x^t - \bar{\alpha}$ is the product of all distinct monic repeated irreducible factors of $\bar{f}(x)$ over $\Z/p\Z$. 
	
	Let $\alpha_1$ and $\alpha_2$ be as defined in Proposition \ref{2.06} and $\alpha = \alpha_1^{u_1} \alpha_2^{u_2} \pmod{p^2}$. Suppose $\beta$ is any repeated root of $\bar f (x)= (x^k+\bar c)^m -\bar a x^n$ in algebraic closure of $\Z/p\Z.$ Then using the derivation similar to those for \eqref{0.009}, we have $\beta^k \equiv \alpha_1 \pmod{p}$. Clearly, $\beta^k \in \Z/p\Z$ and $f(\beta)\equiv 0 \pmod{p}$ imply that
	\begin{align*}
		\beta^n \equiv a^{-1} (\beta^k +c)^m \equiv \alpha_{2} \in {\Z/p\Z}.
	\end{align*}
	As ${\gcd(n,k)}=t$ so there exist $u_1,u_2\in\Z$ such that $ku_1+nu_2=t.$ Using this we obtain $\beta^t = (\beta^k)^{u_1}(\beta^n)^{u_2} \in {\Z/p\Z}$, i.e., $\beta^t \equiv \alpha_1^{u_1} \alpha_{2}^{u_2}\pmod{p}$. Thus, we have proved that any repeated root of $\bar f(x)$ is a root of $ x^t-\bar{\alpha}$. 
	Now using Proposition \ref{2.06}, if $\beta_1$ is a root of $x^t-\bar{\alpha}$  with $\beta_1 \equiv \alpha_1^{u_1} \alpha_{2}^{u_2} \pmod{p}$, then $\beta_1$ is a repeated root of $\bar{f}(x)$ if and only if $p\mid (\calC - \calE)$. Hence, we conclude that every root of $x^t-\bar{\alpha}$ is a repeat root of $f(x)$ modulo $p$. We have shown that every root of $x^t - \bar{\alpha}$ is a repeated root of $f(x)$ modulo $p$.  
%Since $\deg{g(x)}=t\geq 1$, hence $g(x)$ has a root in algebraic closure of $\Z/p\Z$. Let $\beta_0$  be a root of $g(x)$ in algebraic closure of $\Z/p\Z$. 
Since $t=\gcd(n,k)$, it is easy to see that
	\begin{equation}\label{3.20}
		f(x)= ((x^t)^{k_1}+c)^m -a(x^t)^{n_1} = (x^t-\alpha)q(x) +  (\alpha^{k_1}+c)^m - a \alpha^{n_1},
	\end{equation}
	for some $q(x) \in \Z[x^t]$. As $x^t - \bar{\alpha}$ divides $\bar{f}(x)$, we have $\bar{f}(x) = (x^t-\bar{\alpha})\bar{q}(x)$. Let $\bar{f}(x)=\bar g_1(x)^{e_1}\cdots \bar g_t(x)^{e_t}$ be the factorization of $\bar{f}(x)$ into a product of  powers of distinct irreducible polynomials over $\Z /p\Z$ with each $g_{i}(x)\in Z[x]$ monic. If necessary, after renaming assume that $e_{i}>1$ for $1\leq i\leq t_{1}$ and $e_i=1$ for $t_1< i \leq t$. Then  $x^t-\bar{\alpha}=\prod\limits_{i=1}^{t_1}\bar g_i(x)$. Write\\\\
	$x^t-\alpha=\prod\limits_{i=1}^{t_1} g_i(x)+ph_1(x), \quad q(x)=\prod\limits_{i=1}^{t_1}g_i(x)^{e_i-1}~~ \prod\limits_{i=t_1+1}^{t}g_i(x)+ph_2(x)$ \\\\ for some $h_1(x),h_2(x)\in\Z[x]$. Substituting from the above equation in $\eqref{3.20}$, we  have
	\begin{align*}
		f(x)&=\prod_{i=1}^{t}g_i(x)^{e_i}
+ph_1(x)\prod_{i=1}^{t_1}g_i(x)^{e_i-1}	\prod_{i=t_1+1}^{t}g_i(x)+ph_2(x)\prod_{i=1}^{t_1}g_i(x) \nonumber\\ &+~p^2h_1(x)h_2(x) + (\alpha^{k_1}+c)^m - a \alpha^{n_1}.
	\end{align*}
 Clearly, each summand on the right-hand side of the above equation except possibly $(\alpha^{k_1}+c)^m - a \alpha^{n_1}$ belongs to $\langle p, g_i(x) \rangle^2$  for $1\leq i\leq t_1$. So $f(x)$ $\in \langle p, g_i(x) \rangle^2$ for some $i$ $1\leq i\leq t_1$ if and only if $p^2$ divides $(\alpha^{k_1}+c)^m - a \alpha^{n_1}$. As $\bar f(x)$ is not divisible by $\bar g_i(x)^2$ for $t_1< i\leq t$, it is clear that $f(x)\not\in  \langle p, g_i(x) \rangle^2$ for such $i$. So $f(x)\not \in  \langle p, g_i(x) \rangle^2$ for any $i$, $1 \leq i \leq t$ if and only if $p^2\nmid ((\alpha^{k_1}+c)^m - a \alpha^{n_1})$. To prove this case, it only remains to prove that $(\alpha^{k_1}+c)^m - a \alpha^{n_1} \equiv 0 \pmod{p^2}$ if and only if $\calC - \calE \equiv 0\pmod{p^2}$. This follows from Propositions \ref{2.06} and \ref{4.4:prop} and hence the theorem.
\end{proof}

\end{document}
